\section{Các số đặc trưng đo xu thế trung tâm}


\subsection{Đề 1}
\caulc
\Opensolutionfile{ans}[ans/ans-TN-C5-B3-D1]
% Câu 1
\begin{ex}%[0D6N3-5]
	Một cửa hàng trà sữa vừa khai trương, thống kê lượng khách tới quán trong $7$ ngày đầu và thu được mẫu số liệu sau:
	\begin{center}
		\begin{tabular}{|c|c|c|c|c|c|c|}
			\hline
			Ngày $1$ & Ngày $2$ & Ngày $3$ & Ngày $4$ & Ngày $5$ & Ngày $6$ & Ngày $7$\\
			\hline
			$575$ & $454$ & $400$ & $325$ & $351$ & $333$ & $412$\\
			\hline
		\end{tabular}
	\end{center}
	Mệnh đề nào sau đây là đúng?
	\choice
	{Số trung vị là $263$}
	{\True Số trung bình làm tròn đến hàng phần trăm là $407{,}14$}
	{Dấu hiệu điều tra ở đây là doanh thu của quán trà sữa}
	{Ngày $2$ là mốt của mẫu số liệu này}
	\loigiai{
		Ta có $\overline x=\dfrac{575+454+400+325+351+333+412}{7}=407{,}14$.}
\end{ex}
\begin{ex}%[0D6H3-2]
	Điểm trung bình thi học kỳ II môn Toán của một nhóm gồm $N$ học sinh lớp 12A6 là $8{,}1$ . Biết rằng tổng điểm môn toán của nhóm này là $72{,}9$ . Tìm số học sinh của nhóm.
	\choice
	{$20$}
	{\True $9$}
	{$8$}
	{$15$}
	\loigiai{		
		Ta có giá giá trị $N=\dfrac{72{,}9}{8{,}1}=9$ (học sinh).}
\end{ex}
\begin{ex}%[0D6N3-2]
	Tuổi đời của $16$ công nhân trong xưởng sản xuất được thống kê trong bảng sau
	\begin{center}
		\begin{tabular}{|c|c|c|c|c|c|c|c|}
			\hline Tuổi & $25$ & $26$ & $27$ & $29$ & $30$ & $33$ & Cộng \\
			\hline Số người & $2$ & $3$ & $4$ & $3$ & $3$ & $1$ & $16$ \\
			\hline
		\end{tabular}
	\end{center}
	Tìm số trung bình $\overline x $ của mẫu số liệu trên
	\choice
	{$ 28$}
	{$ 27{,}75$}
	{\True $ 27{,}875$}
	{$ 27$}
	\loigiai{		
		Mẫu số liệu có số trung bình là 
		\allowdisplaybreaks
		\begin{eqnarray*}
			\overline x =\dfrac{25\cdot 2+26\cdot 3+27\cdot 4+29\cdot 3+30\cdot 3+33\cdot 1}{16}= 27{,}875.
		\end{eqnarray*}
	}
\end{ex}
\begin{ex}%[0D6N3-3]
	Kết quả điểm kiểm tra môn Toán trong một kì thi của $ 200$ em học sinh được trình bày ở bảng sau:
	\begin{center}
		\begin{tabular}{|c|c|c|c|c|c|c|c|}
			\hline Điềm & $5$ & $6$ & $7$ & $8$ & $9$ & $10$ & Cộng \\
			\hline Tần số & $10$ & $35$ & $38$ & $63$ & $42$ & $12$ & $200$ \\
			\hline
		\end{tabular}
	\end{center}
	Số trung vị của bản phân bố tần số nói trên là
	\choice
	{\True $8$}
	{$ 7$}
	{$ 6$}
	{Đáp án khác}
	\loigiai{
		Số trung vị của bản phân bố tần số nói trên là $8$.}
\end{ex}
\begin{ex}%[0D6N3-3]
	Cho bảng số liệu điểm bài kiểm tra môn toán của $ 20$ học sinh.\\
	\begin{center}
		\begin{tabular}{|c|c|c|c|c|c|c|c|c|}
			\hline Điểm & $4$ & $5$ & $6$ & $7$ & $8$ & $9$ & $10$ & Cộng \\
			\hline Số học sinh & $1$ & $2$ & $3$ & $4$ & $5$ & $4$ & $1$ & $20$ \\
			\hline
		\end{tabular}
	\end{center}
	Tìm số trung vị của bảng số liệu trên.
	\choice
	{$ 8$}
	{$ 7$}
	{$ 7{,}3$}
	{\True $ 7{,}5$}
	\loigiai{
		Số trung vị của bảng số liệu có $20$ số là trung bình cộng của số thứ $10$ và số thứ $11$.\\
		Ta có số thứ $10$ là $7$ ; Số thứ $11$ là $8$.\\
		Do đó $M_e=\dfrac{7+8}{2}=7{,}5$.}
\end{ex}
\begin{ex}%Câu 6
	Cho dãy số liệu thống kê: $48$ , $36$ , $33$ , $38$ , $32$ , $48$ , $42$ , $33$ , $39$. Khi đó số trung vị là
	\choice
	{$32$}
	{$36$}
	{\True $38$}
	{$40$}
	\loigiai{		
		Dãy số liệu thống kê được xếp thành dãy không giảm là $ 32$, $ 33$, $ 33$, $ 36$, $ 38$, $ 39$, $ 42$, $ 48$, $ 48$.\\
		Ta có số trung vị là $M_e=38$.}
\end{ex}
\begin{ex}%[0D6N3-3]
	Cho mẫu thống kê $\left\{28,\,16,\,13,\,18,\,12,\,28,\,13,\,19\right\}$ . Trung vị của mẫu số liệu trên là bao nhiêu?
	\choice
	{$ 14$}
	{\True $ 17$}
	{$ 18$}
	{$ 20$}
	\loigiai{		
		Mẫu thống kê trên có $ 8$ số liệu được sắp xếp theo thứ tự không giảm là $12,\,13,\,13,\,16,\,18,\,19,\,28,\,28$, nên trung vị của mẫu số liệu trên là $M_e=\dfrac{16+18}{2}=17$.}
\end{ex}
\begin{ex}%[0D6N3-3]
	Cho bảng số liệu ghi lại điểm của $ 40$ học sinh trong bài kiểm tra $ 1$ tiết môn Toán.
	\begin{center}
		\begin{tabular}{|c|c|c|c|c|c|c|c|c|c|}
			\hline
			Điểm & $ 3$ & $ 4$ & $ 5$ & $ 6$ & $ 7$ & $ 8$ & $ 9$ & $ 10$ & Cộng\\
			\hline
			Số học sinh & $ 2$ & $ 3$ & $ 7$ & $18$ & $ 3$ & $ 2$ & $ 4$ & $ 1$ & $ 40$\\
			\hline
		\end{tabular}
	\end{center}
	Số trung vị là
	\choice
	{$ 5$}
	{\True $ 6$}
	{$ 6{,}5$}
	{$ 7$}
	\loigiai{
		Số trung vị là $M_e=\dfrac{x_{20}+x_{21}}{2}=\dfrac{6+6}{2}=6$.}
\end{ex}
\begin{ex}%[0D6N3-3]
	Số điểm kiểm tra $11$ môn của một nhóm gồm $11$ học sinh được cho trong bảng sau:
	\begin{center}
		\begin{tabular}{|c|c|c|c|c|c|c|c|}
			\hline
			Điểm & $4$ & $5$ & $7$ & $8$ & $9$ & $10$ &\\
			\hline
			Tần số & $2$ & $1$ & $2$ & $3$ & $1$ & $2$ & N=$11$\\
			\hline
		\end{tabular}
	\end{center}
	Số trung vị của mẫu số liệu trên là
	\choice
	{$7$}
	{\True $7{,}5$}
	{$8$}
	{$8{,}5$}
	\loigiai{
		Số trung vị là $M_e=\dfrac{7+8}{2}=7{,}5$.}
\end{ex}
\begin{ex}%[0D6N3-5]
	Cho mẫu số liêu thống kê $\{6{,}5{,}5{,}2{,}9{,}10{,}8\}$. Mốt của mẫu số liêu trên bằng bao nhiêu?
	\choice
	{\True $5$}
	{$10$}
	{$2$}
	{$6$}
	\loigiai{
		Mốt của mẫu số liêu trên là $M_O=5$.
	}
\end{ex}
\begin{ex}%[0D6N3-3]
	Cho bảng số liệu thống kê chiều cao của một nhóm học sinh như sau:
	\begin{center}
		\begin{tabular}{|l|l|l|l|l|l|l|l|l|l|l|l|l|l|l|l|}
			\hline $150$ & $153$ & $153$ & $154$ & $154$ & $155$ & $160$ & $160$ & $162$ & $162$ & $163$ & $163$ & $163$ & $165$ & $165$ & $167$ \\
			\hline
		\end{tabular}
	\end{center}
	Số trung vị của bảng số liệu nói trên là
	\choice
	{\True $ 161$}
	{$ 153$}
	{$ 163$}
	{$ 156$}
	\loigiai{		
		Ta có trong bảng số liệu thống kê có tất cả $ 16$ giá trị. Do đó số trung vị bằng trung bình cộng của hai số đứng thứ $ 8$ và $ 9$ trong bảng số liệu thống kê.\\
		Ta có $M_e=\dfrac{160+162}{2}=161.$}
\end{ex}
\begin{ex}%[0D6N3-5]
	Cho bảng số liệu nghi lại điểm của $ 40$ học sinh trong bài kiểm tra $ 1$ tiết môn Toán.
	\begin{center}
		\begin{tabular}{|c|c|c|c|c|c|c|c|c|c|}
			\hline
			Điểm & $ 3$ & $ 4$ & $ 5$ & $ 6$ & $ 7$ & $ 8$ & $ 9$ & $ 10$ & Cộng\\
			\hline
			Số học sinh & $ 2$ & $ 3$ & $ 7$ & 18 & $ 3$ & $ 2$ & $ 4$ & $ 1$ & $ 40$\\
			\hline
		\end{tabular}
	\end{center}
	Mốt của dấu hiệu
	\choice
	{$M_0=40$}
	{$M_0=18$}
	{\True $M_0=6$}
	{Không phải các số trên}
	\loigiai{	
		Điểm $ 6$ xuất hiện nhiều nhất nên $M_0=6$.}
\end{ex}
\Closesolutionfile{ans}
\indapan{6}{ans/ans-TN-C5-B3-D1}
%%%%%%%%%%%%%%%%%%%
\cauds
\Opensolutionfile{ans}[ans/ans-TF-C5-B3-D1]
\begin{ex}%[0D6N3-3]
	Tập đoàn $X$ có $24$ công ty. Thống kê cuối năm cho biết doanh thu (đơn vị triệu đồng) của $24$ công ty con như sau:\\
	\centerline{\begin{tabular}{|c|c|c|c|c|c|c|c|}
			\hline
			$35\,432$ & $14\,215$ & $24\,436$ & $13\,978$ & $45\,713$ & $16\,323$ & $37\,488$ & $13\,458$\\
			\hline
			$57\,754$ & $53\,345$ & $80\,234$ & $117\,245$ & $74\,506$ & $86\,851$ & $47\,678$ & $611\,298$\\
			\hline
			$19\,397$ & $48\,644$ & $8\,324$ & $9\,599$ & $94\,338$ & $45\,390$ & $37\,492$ & $811\,854$\\
			\hline
	\end{tabular}}\\
	Xét tính đúng sai của các mệnh đề sau
	\choiceTF
	{Doanh thu thấp nhất là $9\,599$}
	{\True Doanh thu lớn nhất là $811\,854$}
	{\True Số trung bình của mẫu số liệu trên khoảng $100\,208$}
	{\True Số trung vị là $45\,551{,}5$}
	\loigiai{
		Sắp xếp lại mẫu số liệu trên theo thứ tự tăng dần từ trái qua phải, từ trên xuống\\
		\centerline{\begin{tabular}{|c|c|c|c|c|c|c|c|}
				\hline
				$8\,324$ & $9\,599$ & $13\,458$ & $13\,978$ & $14\,215$ & $16\,323$ & $19\,397$ & $24\,436$\\
				\hline
				$35\,432$ & $37\,488$ & $37\,492$ & $45\,390$ & $45\,713$ & $47\,678$ & $48\,644$ & $53\,345$\\
				\hline
				$57\,754$ & $74\,506$ & $80\,234$ & $86\,851$ & $94\,338$ & $117\,245$ & $611\,298$ & $811\,854$\\
				\hline
		\end{tabular}}\\
		Số trung bình của mẫu số liệu trên khoảng $100\,208.$\\
		Số trung vị là $45\,551{,}5$ .\\
		Vì giá trị trung bình lớn hơn giá trị trung vị rất nhiều, điều đó thể hiện mẫu có
		một số giá trị bất thường. Vì vậy, số trung vị làm đại diện mẫu sẽ tốt hơn.
		Do đó
		\begin{itemchoice} 
			\itemch Sai.
			\itemch Đúng. 
			\itemch Đúng. 
			\itemch Đúng.  	
		\end{itemchoice}	
	}
\end{ex}
\begin{ex}%[0D6N3-3]
	Thống kê chiều cao (đơn vị cm) của nhóm $15$ bạn nam lớp $10$ cho kết quả như sau:\\
	\centerline{\begin{tabular}{|c|c|c|c|c|c|c|c|c|c|c|c|c|c|c|}
			\hline
			$162$ & $157$ & $170$ & $165$ & $166$ & $157$ & $159$ & $164$ & $172$ & $155$ & $156$ & $156$ & $180$ & $165$ & $155$\\
			\hline
	\end{tabular}}\\
	Xét tính đúng sai của mệnh đề
	\choiceTF
	{ Chiều cao thấp nhất là $156$}
	{\True $Q_2=162$}
	{ $Q_1=157$}
	{ $Q_3=170$}
	\loigiai{
		Sắp xếp số liệu theo thứ tự tăng dần ta được:\\
		\centerline{\begin{tabular}{|c|c|c|c|c|c|c|c|c|c|c|c|c|c|c|}
				\hline
				$155$ & $155$ & $156$ & $156$ & $157$ & $157$ & $159$ & $162$ & $164$ & $165$ & $165$ & $166$ & $170$ & $172$ & $180$\\
				\hline
		\end{tabular}}\\
		Vì có $15$ giá trị nên số trung vị là số ở vị trí thứ $8$: $ {Q_2}=162$ .\\
		Nửa số liệu bên trái $Q_2$ là\\
		\centerline{\begin{tabular}{|c|c|c|c|c|c|c|}
				\hline
				$155$ & $155$ & $156$ & $156$ & $157$ & $157$ & $159$\\
				\hline
		\end{tabular}}\\
		Ta tìm được trung vị $Q_1=156$ .\\
		Nửa số liệu bên phải $Q_2$ là\\
		\centerline{\begin{tabular}{|c|c|c|c|c|c|c|}
				\hline
				$164$ & $165$ & $165$ & $166$ & $170$ & $172$ & $180$\\
				\hline
		\end{tabular}}\\
		Ta tìm được trung vị $Q_3=166$ .
		Do đó
		\begin{itemchoice} 
			\itemch Sai. 
			\itemch Đúng. 
			\itemch Sai. 
			\itemch Sai. 	
		\end{itemchoice}		
	}
\end{ex}
\begin{ex}%[0D6H3-4]
	Thống kê số bao xi măng được bán ra tại một cửa hàng vật liệu xây dựng trong $24$ tháng cho kết quả như sau:\\
	\centerline{\begin{tabular}{|c|c|c|c|c|c|c|c|}
			\hline
			$72$ & $89$ & $88$ & $73$ & $63$ & $265$ & $69$ & $65$\\
			\hline
			$94$ & $80$ & $81$ & $98$ & $66$ & $71$ & $84$ & $73$\\
			\hline
			$93$ & $59$ & $60$ & $61$ & $83$ & $72$ & $85$ & $66$\\
			\hline
	\end{tabular}}\\
	Xét tính đúng sai của các mệnh để sau
	\choiceTF
	{\True  Mỗi tháng cửa hàng bán trung bình $83{,}75$ bao}
	{ Số trung vị là $72$ }
	{\True  Sai khác giữa số trung bình và số trung vị là $10{,}75$}
	{Khoảng cách từ $Q_1$ đến $Q_2$ là $8$}
	\loigiai{
		Sắp xếp lại mẫu dữ liệu theo thứ tự tăng dần ta được:\\
		\centerline{\begin{tabular}{|c|c|c|c|c|c|c|c|}
				\hline
				$59$ & $60$ & $61$ & $63$ & $65$ & $66$ & $66$ & $69$\\
				\hline
				$71$ & $72$ & $72$ & $73$ & $73$ & $80$ & $81$ & $83$\\
				\hline
				$84$ & $85$ & $88$ & $89$ & $93$ & $94$ & $98$ & $265$\\
				\hline
		\end{tabular}}\\
		Mỗi tháng cửa hàng bán trung bình $83{,}75$ bao.\\
		Số trung vị là $73$ .\\
		Sai khác giữa số trung bình và số trung vị là $10{,}75$ . Điều này nói lên rằng trong\\
		mẫu có một số giá trị bất thường.\\
		Ta có số trung vị $Q_2=73$ .\\
		Số trung vị của nửa bên trái $Q_2$ là $Q_1=66$ .\\
		Số trung vị nửa bên phải $Q_2$ là $Q_3=\dfrac{85+88}{2}=86{,}5$ .\\
		Khoảng cách từ $Q_1$ đến $Q_2$ là $7$ , từ $Q_2$ đến $Q_3$ là $13{,}5$ .\\
		Điều này nói lên rằng mẫu số liệu tập trung với mật độ cao ở bên trái của $Q_2$.
		Do đó
		\begin{itemchoice} 
			\itemch Đúng. 
			\itemch Sai. 
			\itemch Đúng. 
			\itemch Sai. 	
		\end{itemchoice}
	}
\end{ex}
\begin{ex}%[0D6H3-5]
	Cho mẫu số liệu thống kê về sản lượng chè thu được trong $1$ năm (kg/sào) của $10$ hộ gia đình:\\
	\centerline{\begin{tabular}{|c|c|c|c|c|c|c|c|c|c|}
			\hline
			$112$ & $111$ & $112$ & $113$ & $114$ & $116$ & $115$ & $114$ & $115$ & $114$\\
			\hline
	\end{tabular}}\\
	Xét tính đúng sai của các mệnh đề sau
	\choiceTF
	{\True  Sản lượng chè trung bình thu được trong một năm của mỗi gia đình là $\approx 113{,}6$ (kg/sào)}
	{Ta viết lại mẫu số liệu trên theo thứ tự không giảm là	$111$, $112$, $112$, $113$, $114$, $114$, $114$, $115$, $115$, $116$}
	{\True Số trung vị là $113$}
	{\True  $114$ là mốt của mẫu số liệu đã cho}
	\loigiai{
		Sản lượng chè trung bình thu được trong một năm của mỗi gia đình là \\$\overline x=\dfrac{112+111+112+113+114+116+115+114+115+114}{10}\approx 113{,}6$ (kg/sào).\\
		Ta viết lại mẫu số liệu trên theo thứ tự không giảm:
		\allowdisplaybreaks
		\begin{eqnarray*}
			\begin{array}{*{20}{l}}
				{111}&{112}&{112}&{113}&{114}&{114}&{114}&{115}&{115}&{116}
			\end{array}
		\end{eqnarray*}
		Vì số giá trị của mẫu $n=10$ (chẵn) nên trung bình cộng hai số chính giữa mẫu chính là trung vị, vậy trung vị là $\dfrac{114+114}{2}=114$.
		\begin{itemchoice} 
			\itemch Đúng. 
			\itemch Sai. 
			\itemch Đúng. 
			\itemch Đúng. 	
		\end{itemchoice}
	}
\end{ex}
\Closesolutionfile{ans}
\indapan{3}{ans/ans-TF-C5-B3-D1}
\Opensolutionfile{ans}[ans/ans-KQ-C5-B3-D1]
\caukq
\begin{ex}%[0D6H3-2]
	Một nhóm $11$ học sinh tham gia một kì thi. Số điểm thi của $11$ học sinh đó được sắp xếp từ thấp đến cao theo thang điểm $100$ như sau: $ 0;0;63$; $ 65;69;70;72;78;81;85;89$. Tìm điểm số trung bình của nhóm $11$ học sinh này?\shortans[]{$61{,}1$}
	\loigiai{
		Điểm trung bình là $\overline x=\dfrac{0+0+63+\cdots+85+89}{11}=\dfrac{672}{11}\approx 61{,}1$.}
\end{ex}
\begin{ex}%[0D6H3-2]
	Một cửa hàng bán $6$ loại quạt điện với giá tiền là  $100{,}1$ $50$, $300$, $350$, $400$, $500$ (nghìn đồng). Số quạt điện mà cửa hàng bán ra trong mùa hè vừa qua được thống kê trong bảng tần số sau:\\
	\centerline{\begin{tabular}{|c|c|c|c|c|c|c|}
			\hline
			Giá tiền & $100$ & $150$ & $300$ & $350$ & $400$ & $500$\\
			\hline
			Số quạt bán được & $15$ & $25$ & $31$ & $26$ & $12$ & $4$\\
			\hline
	\end{tabular}}\\
	Tìm số tiền trung bình mà cửa hàng thu được khi bán mỗi chiếc quạt?\shortans[]{$269$}
	\loigiai{
		$\overline x=\dfrac{100\cdot 15+150\cdot 25+300\cdot 31+350\cdot 26+400\cdot 12+500\cdot 4}{15+25+31+26+12+4}\approx 269$ (nghìn).}
\end{ex}
\begin{ex}%[0D6H2-1]
	Trong $7$ tháng đầu năm, số sản phẩm sản xuất mỗi tháng của công ty X đều tăng trưởng khoảng $ 5\% $ so với tháng trước đó. Biết rằng, trong bảng dưới đây, số sản phẩm sản xuất của một tháng bị nhập sai, Hãy tìm tháng đó.\\
	\centerline{\begin{tabular}{|c|c|c|c|c|c|c|c|}
			\hline
			Tháng & $1$ & $2$ & $3$ & $4$ & $5$ & $6$ & $7$\\
			\hline
			Số sản phẩm sản xuất & $500$ & $525$ & $551$ & $569$ & $606$ & $636$ & $668$\\
			\hline
	\end{tabular}}\shortans[]{$4$}
	\loigiai{
		Tỉ lệ phần trăm tăng thêm của số sản phẩm bán ra mỗi tháng được tính ở bảng dưới đây.\\
		\centerline{\begin{tabular}{|c|c|c|c|c|c|c|}
				\hline
				Tháng & $2$ & $3$ & $4$ & $5$ & $6$ & $7$\\
				\hline
				Tỉ lệ phần trăm tăng thêm so với tháng trước & $ 5\% $ & $ 4{,}95\% $ & $ 3{,}4\% $ & $ 6{,}5\% $ & $ 4{,}95\% $ & $ 5{,}03\% $\\
				\hline
		\end{tabular}}\\
		Ta thấy tỉ lệ tăng của tháng $4$ và tháng $5$ đều khác xa $ 5\% $.\\
		Do đó trong bảng số liệu đã cho, số sản phẩm của tháng $4$ là không chính xác.}
\end{ex}
\begin{ex}%[0D6H3-4]
	Hãy tìm các tứ phân vị của mẫu số liệu sau: ${2;\, 3; \, 7; \, 4; \, 5; \, 3; \, 4; \, 4}$. Khi đó $Q_1+Q_2+Q_3$ có giá trị bằng bao nhiêu?
	\shortans{$11{,}5$}
	\loigiai{
		Sắp xếp mẫu: ${2;\, 3; \, 7; \, 4; \, 5; \, 3; \, 4; \, 4}$. Kích thước mẫu là $8$ (chẵn).\\
		Khi đó trung vị của mẫu là trung bình cộng của giá trị thứ $4$ và $5$ là ${M_e=\dfrac{1}{2}(4+4)=4}$.
		Tứ phân vị thứ nhất là ${Q_1=\dfrac{1}{2}(3+3)=3}$.
		Tứ phân vị thứ hai là ${Q_2=M_e=4}$.
		Tứ phân vị thứ ba là ${Q_3=\dfrac{1}{2}(4+5)=4,5}$.\\
		Vậy $Q_1+Q_2+Q_3=11{,}5$.
	}
\end{ex}
\begin{ex}%[0D6H2-2]
	Bảng sau thống kê số lớp và số học sinh theo từng khối ở một trường Trung học phổ thông.\\
	\centerline{\begin{tabular}{|c|c|c|c|}
			\hline
			Khối & $10$ & $11$ & $12$\\
			\hline
			Số lớp & $14$ & $13$ & $15$\\
			\hline
			Số học sinh & $555$ & $519$ & $615$\\
			\hline
	\end{tabular}}\\
	Hiệu trưởng trường đó cho biết sĩ số mỗi lớp trong trường đều không vượt quá $40$ học sinh. Biết rằng trong bảng trên có một khối lớp bị thống kê sai, hãy tìm khối lớp đó.
	\shortans[]{$12$}
	\loigiai{
		Theo bảng thống kê đã cho, sĩ số tối đa của mỗi lớp theo từng khối cho ở bản sau:\\
		\centerline{\begin{tabular}{|c|c|c|c|}
				\hline
				Khối & $10$ & $11$ & $12$\\
				\hline
				Sĩ số tối đa mỗi khối & $560$ & $520$ & $600$\\
				\hline
		\end{tabular}}\\
		Theo thông tin hiệu trưởng cung cấp thì thông tin Khối $12$ đã bị thống kê sai vì Hiệu trưởng trường đó cho biết sĩ số mỗi lớp trong trường đều không vượt quá $40$ học sinh nhưng khi thống kê thì sĩ số khối $12$ vượt quá số học sinh tối đa ($600$ học sinh).}
\end{ex}
\begin{ex}%[0D6H2-2]
	Trong giờ học Toán lớp $7$, giáo viên giao nhiệm vụ cho mỗi nhóm đo các góc của một tam giác. Kết quả được ghi lại trong bảng sau:\\
	\centerline{\begin{tabular}{|c|c|c|c|c|}
			\hline
			Nhóm & Nhóm $1$ & Nhóm $2$ & Nhóm $3$ & Nhóm $4$\\
			\hline
			Góc thứ nhất & $35^\circ$ & $61^\circ$ & $33^\circ$ & $100^\circ$\\
			\hline
			Góc thứ hai & $77^\circ$ & $74^\circ$ & $102^\circ$ & $37^\circ$\\
			\hline
			Góc thứ ba & $68^\circ$ & $45^\circ$ & $47^\circ$ & $43^\circ$\\
			\hline
	\end{tabular}}\\
	Trong bảng trên có nhóm ghi kết quả sai. Hãy tìm nhóm nào sai?
	\shortans{$3$}
	\loigiai{
		Do tổng ba góc trong một tam giác bằng $180^\circ$ nên ta lập bảng sau:\\
		\centerline{\begin{tabular}{|c|c|c|c|c|}
				\hline
				Nhóm & Nhóm $1$ & Nhóm $2$ & Nhóm $3$ & Nhóm $4$\\
				\hline
				Góc thứ nhất & $35^\circ$ & $61^\circ$ & $33^\circ$ & $100^\circ$\\
				\hline
				Góc thứ hai & $77^\circ$ & $74^\circ$ & $102^\circ$ & $37^\circ$\\
				\hline
				Góc thứ ba & $68^\circ$ & $45^\circ$ & $47^\circ$ & $43^\circ$\\
				\hline
				Tổng ba góc & $180^\circ$ & $180^\circ$ & $182^\circ$ & $180^\circ$\\
				\hline
		\end{tabular}}\\
		Vậy nhóm $3$ đo sai vì tổng số đo của ba góc khác $180^\circ$ .}
\end{ex}
\Closesolutionfile{ans}
\indapan{6}{ans/ans-KQ-C5-B3-D1}
\newpage
\subsection{Đề 2}
\caulc
\Opensolutionfile{ans}[ans/ans-TN-C5-B3-D2]
\begin{ex}%[0D6N3-5]
	Tiền thưởng (triệu đồng) của cán bộ và nhân viên trong một công ty được cho ở bảng dưới đây:\\
	\begin{center}
		\begin{tabular}{|c|c|c|c|c|c|c|}
			\hline Tiền thưởng & $12$ & $13$ & $14$ & $15$ & $16$ & Cộng \\
			\hline Tần số & $25$ & $15$ & $11$ & $16$ & $17$ & $84$ \\
			\hline
		\end{tabular}
	\end{center}
	Tính mốt $M_o$.
	\choice
	{$M_o=15$}
	{\True $M_o=12$}
	{$M_o=10$}
	{$M_o=16$}
	\loigiai{
		Vì tần số bằng $25$ là lớn nhất và ứng với số tiền thưởng (triệu đồng) của cán bộ và nhân viên trong một công ty là $12$ nên $M_o=12$.}
\end{ex}

\begin{ex}%[0D6N3-5]
	Cho bảng phân bố tần số tiền thưởng (triệu đồng) cho cán bộ và nhân viên trong một công ty\\
	\centerline{\begin{tabular}{|c|c|c|c|c|c|c|}
			\hline
			Tiền thưởng & $2$ & $3$ & $4$ & $5$ & $6$ & Cộng\\
			\hline
			Tần số & $5$ & $15$ & $10$ & $6$ & $7$ & $43$\\
			\hline
	\end{tabular}}\\
	Mốt của bảng phân bố tần số đã cho là
	\choice
	{\True $3$ triệu đồng}
	{$2$ triệu đồng}
	{$6$ triệu đồng}
	{$5$ triệu đồng}
	\loigiai{
		Mốt của bảng phân bố tần số là giá trị (xi) có tần số (ni) lớn nhất và được kí hiệu là $M_o$}
\end{ex}

\begin{ex}%[0D6N3-5]
	Kết quả thi môn Toán giữa kì 1$1$ của lớp $10A_3$ trường THPT Ba Vì được thống kê như sau:\\
	\begin{center}
		\begin{tabular}{|c|c|c|c|c|c|c|c|}
			\hline Điểm thi & $5$ & $6$ & $7$ & $8$ & $9$ & $10$ & Cộng \\
			\hline Tần số & $5$ & $7$ & $8$ & $12$ & $8$ & $5$ & $45$ \\
			\hline
		\end{tabular}
	\end{center}
	Giá trị mốt $M_o$ của bảng phân bố tần số trên bằng
	\choice
	{$5$}
	{$7$}
	{\True $8$}
	{$12$}
	\loigiai{
		Mốt của bảng phân bố tần suất là giá trị có tần số lớn nhất nên ta có $M_o=8$.}
\end{ex}

\begin{ex}%[0D6N3-5]
	Mốt của một bảng phân bố tần số là
	\choice
	{tần số lớn nhất trong bảng phân bố tần số}
	{\True giá trị có tần số lớn nhất trong bảng phân bố tần số}
	{giá trị có tần số nhỏ nhất trong bảng phân bố tần số}
	{tần số nhỏ nhất trong bảng phân bố tần số}
	\loigiai{
		Mốt của một bảng phân bố tần số là giá trị có tần số lớn nhất.}
\end{ex}

\begin{ex}%[0D6N3-2]
	Điểm số của $100$ học sinh tham dự kỳ thi học sinh giỏi toán ở tỉnh $A$ (thang điểm là $20$) được thống kê theo bảng sau:\\
	\centerline{\begin{tabular}{|c|c|c|c|c|c|c|c|c|c|c|c|}
			\hline
			Điểm $(x)$ & 9 & 10 & 11 & 12 & 13 & 14 & 15 & 16 & 17 & 18 & 19\\
			\hline
			Tần số $(n)$ & 1 & 1 & 3 & 5 & 8 & 13 & 19 & 24 & 14 & 10 & 2\\
			\hline
	\end{tabular}}\\
	Trung bình cộng của bảng số liệu trên là
	\choice
	{$15$}
	{$15{,}50$}
	{$16$}
	{\True $15{,}23$}
	\loigiai{
		
		Trung bình cộng của bảng số liệu trên là
		\allowdisplaybreaks
		\begin{eqnarray*}
			\dfrac{9\cdot 1+10\cdot 1+11\cdot 3+12\cdot 5+13\cdot 8+14\cdot 13+15\cdot 19+16\cdot 24+17\cdot 14+18\cdot 10+19\cdot 2}{100}=15{,}23.
	\end{eqnarray*}}
\end{ex}

\begin{ex}%[0D6N3-2]
	Điều tra về số tiền mua đồ dùng học tập trong một tháng của 40 học sinh, ta có mẫu số liệu như sau (đơn vị: nghìn đồng):\\
	\begin{center}
		\begin{tabular}{|l|l|l|l|l|l|l|l|}
			\hline Giá trị $(x)$ & {$[10 ; 15)$} & {$[15 ; 20)$} & {$[20 ; 25)$} & {$[25 ; 30)$} & {$[30 ; 35)$} & {$[35 ; 40)$} & Cộng \\
			\hline Tần số $(n)$ & $2$ & $5$ & $15$ & $8$ & $9$ & $1$ & $\mathrm{~N}=40$ \\
			\hline
		\end{tabular}
	\end{center}
	Số trung bình của mẫu số liệu là
	\choice
	{$22{,}5$}
	{\True $25$}
	{$25{,}5$}
	{$27$}
	\loigiai{
		Số trung bình cộng của mẫu số liệu là\\
		\allowdisplaybreaks
		\begin{eqnarray*}
			\overline x=\dfrac{12{,}5\cdot 2+17{,}5\cdot 5+22{,}5\cdot 15+27{,}5\cdot 8+32{,}5\cdot 9+37{,}5\cdot 1}{40}=25.
	\end{eqnarray*}}
\end{ex}

\begin{ex}%[0D6N3-2]
	Tuổi đời của $16$ công nhân trong xưởng sản xuất được thống kê trong bảng sau\\
	\begin{center}
		\begin{tabular}{|l|l|l|l|l|l|l|l|}
			\hline Tuổi & $25$ & $26$ & $27$ & $29$ & $30$ & $31$ & Cộng \\
			\hline Số người & $2$ & $3$ & $4$ & $3$ & $3$ & $1$ & $16$ \\
			\hline
		\end{tabular}
	\end{center}
	Số trung bình $\overline x $ của mẫu số liệu trên là
	\choice
	{$28$}
	{\True $27{,}75$}
	{$27{,}875$}
	{$27$}
	\loigiai{
		Số trung bình $\overline x $ của mẫu số liệu trên là
		\allowdisplaybreaks
		\begin{eqnarray*}
			\dfrac{25\cdot 2+26\cdot 3+27\cdot 4+29\cdot 3+30\cdot 3+31\cdot 1}{16}=27{,}75.
	\end{eqnarray*}}
\end{ex}

\begin{ex}%[0D6N3-2]
	Ba nhóm học sinh gồm $10$ người, $15$ người, $25$ người. Cân nặng trung bình của mỗi nhóm lần lượt là $50$ kg, $38$ kg, $40$ kg. Khối lượng trung bình của ba nhóm học sinh đó là
	\choice
	{$26$ kg}
	{\True $41{,}4$ kg}
	{$42{,}4$ kg}
	{$37$ kg}
	\loigiai{
		
		Khối lượng trung bình của ba nhóm học sinh là\\
		\allowdisplaybreaks
		\begin{eqnarray*}
			\dfrac{1}{10+15+25}\cdot \left(10\cdot 50+15\cdot 38+25\cdot 40\right)=41{,}4\, \text{kg}.
	\end{eqnarray*}}
\end{ex}

\begin{ex}%[0D6N3-2]
	Kết quả điểm kiểm tra 45 phút môn Hóa Học của 100 em học sinh được trình bày ở bảng sau:\\
	\begin{center}
		\begin{tabular}{|l|c|c|c|c|c|c|c|c|c|}
			\hline Điểm & $3$ & $4$ & $5$ & $6$ & $7$ & $8$ & $9$ & $10$ & Cộng \\
			\hline Tần số & $3$ & $5$ & $14$ & $14$ & $30$ & $22$ & $7$ & $5$ & $100$ \\
			\hline
		\end{tabular}
	\end{center}
	Số trung bình cộng của bảng phân bố tần số nói trên là
	\choice
	{\True $6{,}82$}
	{$4$}
	{$6{,}5$}
	{$7{,}22$}
	\loigiai{
		Số trung bình cộng của bảng phân bố tần số nói trên là
		\allowdisplaybreaks
		\begin{eqnarray*}
			\overline x=\dfrac{3\cdot 3+4\cdot 5+5\cdot 14+6\cdot 14+7\cdot 30+8\cdot 22+9\cdot 7+10\cdot 5}{100}=6{,}82.
	\end{eqnarray*}}
\end{ex}

\begin{ex}%[0D6N3-2]
	Cho bảng phân bố tần số sau: khối lượng $20$ học sinh lớp $10A$\\
	\begin{center}
		\begin{tabular}{|c|c|}
			\hline Khối lượng (kg) & Tần số \\
			\hline $50$ & $4$ \\
			\hline $51$ & $5$ \\
			\hline $52$ & $6$ \\
			\hline $55$ & $3$ \\
			\hline $56$ & $2$ \\
			\hline
		\end{tabular}
	\end{center}
	Số trung bình cộng $\overline x $ của bảng số liệu đã cho là
	\choice
	{$\overline x=53$}
	{$\overline x=52{,}8$}
	{\True $\overline x=52{,}2$}
	{$\overline x=52$}
	\loigiai{
		Giá trị trung bình $\overline x=\dfrac{50\cdot 4+51\cdot 5+52\cdot 6+55\cdot 3+56\cdot 2}{20}=52{,}2$.}
\end{ex}

\begin{ex}%[0D6N3-3]
	Điểm thi học kì của một học sinh như sau: $4;\,6;\,2;\,7;\,3;\,5;\,9;\,8;\,7;\,10;\,9$. Số trung bình và số trung vị lần lượt là
	\choice
	{$6{,}22$ và $7$}
	{$7$ và $6$}
	{\True $6{,}36$ và $7$}
	{$6$ và $6$}
	\loigiai{
		Mẫu thống kê trên có $11$ số liệu được sắp xếp theo thứ tự không giảm là $2;\,3;\,4;\,5;\,6;\,7;\,7;\,8;\,9;\,9;\,10$.\\
		Số trung bình là $\overline x=\dfrac{2+3+4+5+6+2\cdot 7+8+2\cdot 9+10}{11}$ $=\dfrac{70}{11}\approx 6{,}36$.\\
		Trung vị của mẫu số liệu trên là $M_e=7$.
	}
\end{ex}

\begin{ex}%[0D6N3-5]
	Cho bảng phân bố tần số như sau:\\
	\begin{center}
		\begin{tabular}{|c|c|c|c|c|c|c|c|c|}
			\hline Giá trị & $x_1$ & $x_2$ & $x_3$ & $x_4$ & $x_5$ & $x_6$ & $x_7$ & $x_8$ \\
			\hline Tần số & $15$ & $9 n-1$ & $12$ & $n^2+7$ & $14$ & $10$ & $9 n-20$ & $17$ \\
			\hline
		\end{tabular}
	\end{center}
	Tìm $n$ để $M_{_o}^{(1)}=x_2\,;\,M_{_o}^{(2)}=x_4$ là hai mốt của bảng số liệu trên.
	\choice
	{$n=1\,;\,n=8$}
	{\True $n=8$}
	{$n=1$}
	{$n=9$}
	\loigiai{
		Ta có $M_{_o}^{(1)}=x_2\,;\,M_{_o}^{(2)}=x_4$ là hai mốt của bảng phân bố tần số nên 
		\allowdisplaybreaks
		\begin{eqnarray*}
			\heva{&n^2+7=9n-1\\
				&9n-1 > 17}
			\Leftrightarrow\heva{&n^2-9n+8=0\\
				&n > 2}\Leftrightarrow\heva{&\hoac{&n=1\,(\text{loại})\\
					&n=8\,(\text{thỏa mãn})}\\
				&n > 2}\Rightarrow n=8.
		\end{eqnarray*}
	}
\end{ex}
\Closesolutionfile{ans}
\indapan{6}{ans/ans-TN-C5-B3-D2}
%%%%%%%%%%%%%%%%%%%
\cauds
\Opensolutionfile{ans}[ans/ans-TF-C5-B3-D2]
\begin{ex}%[0D6N3-3]
	Bảng số liệu sau cho biết mức lương hàng năm của các cán bộ và nhân viên trong một công ty (đơn vị: nghìn đồng).\\
	\centerline{\begin{tabular}{|c|c|c|c|}
			\hline
			$20910$ & $76000$ & $20350$ & $20060$\\
			\hline
			$21410$ & $20110$ & $21410$ & $21360$\\
			\hline
			$20350$ & $21130$ & $20960$ & $125000$\\
			\hline
	\end{tabular}}\\
	Xét tính đúng sai của các mệnh đề sau
	\choiceTF
	{\True Mức lương trung bình các cán bộ nhân viên là $\overline x=34087{,}5$ (nghìn đồng)}
	{Mức lương lớn nhất là $ 76000$}
	{\True Số trung vị là $ 21045$ (nghìn đồng)}
	{Có thể lấy mức lương bình quân làm giá trị đại diện}
	\loigiai{
		Mức lương trung bình các cán bộ nhân viên là\\
		$\overline x=\dfrac{20910+76000+\ldots+125000}{12}=\dfrac{68175}{2}=34087{,}5$ (nghìn đồng).\\
		Sắp theo thứ tự không giảm bảng lương ta được: $20060\;\;20110\;\;20350$ $ 20350\,\,\,20910\,\,\,20960\,\,\,21130\,\,\,21360\,\,\,21410\,\,\,21410\,\,\,76000\,\,\,125000;\quad (n=12)$.\\
		Số trung vị là $\dfrac{20960+21130}{2}=21045$ (nghìn đồng).\\
		Ý nghĩa: Số trung vị đại diện cho mức lương trung bình của cán bộ nhân viên của công ty, vì trong trường hợp này chênh lệch giữa các số liệu quá lớn nên không thể lấy mức lương bình quân làm giá trị đại diện.\\
		Do đó
		\begin{itemchoice} 
			\itemch Đúng. 
			\itemch Sai. 
			\itemch Đúng. 
			\itemch Sai. 	
		\end{itemchoice}
	}
\end{ex}

\begin{ex}%[0D6N3-4]
	Điểm số của một nhóm $14$ em học sinh sau khi tham gia lớp phụ đạo môn toán được cho ở mẫu sau: $\begin{array}{*{20}{l}}
		3&3&4&4&5&5&5&6&6&6&7&7&8&{10.}
	\end{array}$ Xét tính đúng sai của mệnh đề
	\choiceTF
	{\True $ n=14$}
	{\True $Q_2=5{,}5$}
	{$Q_1=5$}
	{\True $Q_3=7$}
	\loigiai{
		Xét mẫu đã cho: $ 3\,\,\,3\,\,\,4\,\,\,4\,\,\,5\,\,\,5\,\,\,5\,\,\,6\,\,\,6\,\,\,6\,\,\,7\,\,\,7\,\,\,8\,\,\,10\,\left(n=14\right)$.\\
		Tứ phân vị thứ hai là $Q_2=\dfrac{5+6}{2}=5{,}5$.\\
		Xét nửa mẫu bên trái: $ 3\,\,\,3\,\,\,4\,\,\,4\,\,\,5\,\,\,5\,\,\,5$.\\
		Tứ phân vị thứ nhất là $Q_1=4$\\
		Xét nửa mẫu bên phải: $\begin{array}{*{20}{l}}
			6&6&6&7&7&8&{10}
		\end{array}$.\\
		Tứ phân vị thứ ba là $Q_3=7$.\\
		Vậy tứ phân vị của mẫu gồm: $Q_1=4$, $Q_2=5{,}5$, $Q_3=7$.\\
		Do đó
		\begin{itemchoice} 
			\itemch Đúng. 
			\itemch Sai. 
			\itemch Đúng. 
			\itemch Đúng. 	
		\end{itemchoice}
		
	}
\end{ex}

\begin{ex}%[0D6N3-5]
	Cho mẫu số liệu sau: $4$; $5$; $6$; $7$; $8$; $4$; $9$; $4$; $3$. Xét tính đúng sai của mệnh đề
	\choiceTF
	{\True Số trung bình: $\overline x=5{,}5$}
	{Mốt: $M_o=3$}
	{Trung vị là $M_e=4$}
	{\True Tứ phân vị thứ ba là $Q_3=7$}
	\loigiai{
		Số trung bình: $\overline x=\dfrac{55}{10}=5{,}5$. Sắp xếp mẫu: $ 3;4;4;4;5;5;6;7;8;9$. Kích thước mẫu là 10 (chẵn) nên trung vị là $M_e=\dfrac{1}{2}(5+5)=5$. Tứ phân vị thứ nhất là $Q_1=4$. Tứ phân vị thứ hai là $Q_2=M_e=5$. Tứ phân vị thứ ba là $Q_3=7$.\\
		Mốt $M_o=4$.\\
		Do đó 
		\begin{itemchoice} 
			\itemch Đúng. 
			\itemch Sai. 
			\itemch Sai. 
			\itemch Đúng. 	
		\end{itemchoice}
	}
\end{ex}

\begin{ex}%[0D6N3-5]
	Cho mẫu số liệu sau: $6$; $5$; $6$; $7$; $8$; $4$; $6$; $5$; $4$. Xét tính đúng sai của mệnh đề
	\choiceTF
	{\True Số trung bình: $\overline x\approx 5{,}67$}
	{\True Trung vị là $M_e=6$}
	{Tứ phân vị thứ nhất là $Q_1=4$}
	{Mốt: $M_o=5$}	
	\loigiai{
		Số trung bình: $\overline x=\dfrac{51}{9}\approx 5{,}67$.\\
		Sắp xếp mẫu: $4$; $4$; $5$; $5$; $6$; $6$; $6$; $7$; $8$.\\
		Kích thước mẫu là $9$ (lẻ) nên trung vị là $M_e=6$.\\
		Tứ phân vị thứ nhất là $Q_1=\dfrac{1}{2}(4+5)=4{,}5$.\\
		Tứ phân vị thứ hai là $Q_2=M_e=6$.\\
		Tứ phân vị thứ ba là $Q_3=\dfrac{1}{2}(6+7)=6{,}5$.\\
		Mốt: $M_o=6$.\\
		Do đó 
		\begin{itemchoice} 
			\itemch Đúng. 
			\itemch Đúng. 
			\itemch Sai. 
			\itemch Sai. 	
		\end{itemchoice}
	}
\end{ex}
\Closesolutionfile{ans}
\indapan{3}{ans/ans-TF-C5-B3-D2}
\Opensolutionfile{ans}[ans/ans-KQ-C5-B3-D2]
\caukq
\begin{ex}%[0D6N3-2]
	Một công ty vận chuyển $A$ dự kiến thưởng cho nhân viên giao hàng $B$ vào cuối năm dựa vào số đơn hàng giao được trong năm. Số đơn hàng của nhân viên ${{B}}$ giao được trong các tháng được cho trong dãy sau:
	\allowdisplaybreaks
	\begin{eqnarray*}
		{\begin{array}{llllllllllll}
				1002 & 510 & 430 & 395 & 400 & 401 & 396 & 299 & 450 & 450 & 560 & 611.
			\end{array}
		}
	\end{eqnarray*}
	Tính số đơn hàng trung bình giao được trong $1$ tháng của nhân viên ${B}$
	\shortans[]{$439$} 
	\loigiai{
		Số đơn hàng trung bình giao được trong $1$ tháng của nhân viên ${B}$ :
		\allowdisplaybreaks
		\begin{eqnarray*}
			\bar{x} =\dfrac{1002+510+430+395+400+401+396+299+450+450+560+611}{12}=439 \text{(đơn hàng)}.
		\end{eqnarray*}
	}
\end{ex}
\begin{ex}%[0D6N3-3]
	Chiều dài (đơn vị feet) của $7$ con cá voi trưởng thành được cho như sau:
	\allowdisplaybreaks
	\begin{eqnarray*}
		\begin{array}{lllllll}
			48 & 53 & 51 & 31 & 53 & 112 & 52.
		\end{array}
	\end{eqnarray*}	Tìm trung vị của mẫu số liệu trên.
	\shortans[]{$52$}
	\loigiai{
		Sắp xếp lại mẫu số liệu không giảm như sau:
		${\begin{array}{lllllll}
				31 & 48 & 51 & 52 & 53 & 53 & 112
			\end{array}
		}$\\
		Do đó trung vị của mẫu là ${M_e=52}$.		
	}
\end{ex}
\begin{ex}%[0D6N3-2]
	Cho mẫu số liệu có bảng tần suất như sau:\\
	\centerline{\begin{tabular}{|c|c|c|c|c|c|}
			\hline
			Giá trị $x_i$ & $4$ & $5$ & $6$ & $7$ & $8$\\
			\hline
			Tần số tương đối $f_i$ & $0{,}1$ & $0{,}45$ & $0{,}2$ & $0{,}1$ & $0{,}15$\\
			\hline
	\end{tabular}}\\
	Ta có số trung bình của mẫu số liệu là
	\shortans[]{$5{,}75$}
	\loigiai{
		Ta có số trung bình của mẫu số liệu là
		\allowdisplaybreaks
		\begin{eqnarray*}
			\overline  x=4\cdot 0{,}1+5\cdot 0{,}45+6\cdot 0{,}2+7\cdot 0{,}1+8\cdot 0{,}15=5{,}75.
	\end{eqnarray*}}
\end{ex}

\begin{ex}%[0D6N3-4]
	Hãy tìm các tứ phân vị của mẫu số liệu sau: $2$; $3$; $7$; $4$; $5$; $3$; $4$; $4$. Khi đó $Q_1+Q_3$ có giá trị bằng bao nhiêu?
	\shortans[]{$7{,}5$}
	\loigiai{
		Sắp xếp mẫu: $2$; $3$; $3$; $4$; $4$; $4$; $5$; $7$. Kích thước mẫu là $8$ (chẵn).\\
		Khi đó trung vị của mẫu là trung bình cộng của giá trị thứ $4$ và $5$ là $M_e=\dfrac{1}{2}(4+4)=4$.\\
		Tứ phân vị thứ nhất là $Q_1=\dfrac{1}{2}(3+3)=3$.\\
		Tứ phân vị thứ hai là $Q_2=M_e=4$.\\
		Tứ phân vị thứ ba là $Q_3=\dfrac{1}{2}(4+5)=4{,}5$.\\
		Vậy $Q_{1}+Q_{3}=7{,}5$.
	}
\end{ex}

\begin{ex}%[0D6N2-1]
	Trong $6$ tháng đầu năm, số sản phẩm bán ra mỗi tháng của một cửa hàng đều tăng khoảng $25\% $ so với tháng trước đó. Biết rằng, trong bảng dưới đây, số sản phẩm bán ra của một tháng bị nhập sai. Hãy tìm tháng đó.\\
	\centerline{\begin{tabular}{|c|c|c|c|c|c|c|}
			\hline
			Tháng & $1$ & $2$ & $3$ & $4$ & $5$ & $6$\\
			\hline
			Số sản phẩm bán ra & $145$ & $180$ & $225$ & $279$ & $390$ & $435$\\
			\hline
	\end{tabular}}\\
	\shortans[]{$5$}
	\loigiai{
		Tỉ lệ phần trăm tăng thêm của số sản phẩm bán ra mỗi tháng được tính ở bảng dưới đây\\
		\centerline{\begin{tabular}{|c|c|c|c|c|c|}
				\hline
				Tháng & $2$ & $3$ & $4$ & $5$ & $6$\\
				\hline
				Tỉ lệ phần trăm tăng thêm so với tháng trước & $24{,}1\% $ & $25\% $ & $24\% $ & $39{,}8\% $ & $11{,}5\% $\\
				\hline
		\end{tabular}}\\
		Tỉ lệ tăng của tháng $5$ và tháng $6$ đều rất khác so với $25\% $, do đó số liệu trong tháng 5 là dữ liệu đã nhập sai.}
\end{ex}
%
\begin{ex}%[0D6H3-4]
	Bảng sau đây cho biết số lần học tiếng Anh trên internet trong một tuần của một học sinh lớp $10$:
	\begin{center}
			\begin{longtable}{|l|l|l|l|l|l|l|}
			\hline
			Số lần & $0$ & $1$ & $2$ & $3$ & $4$ & $5$ \\
			\hline
			Số học sinh & $2$ & $4$ & $6$ & $12$ & $8$ & $3$ \\
			\hline
		\end{longtable}
	\end{center}
	Hãy tìm các tứ phân vị cho mẫu số liệu này. Tính $T=Q_{1}+Q_{2}+Q_{3}$.\shortans[]{$9$}
	\loigiai{
		Cỡ mẫu: $n=2+4+6+12+8+3=35$.\\
		Vì cỡ mẫu là ${n=35}$, là số lẻ, nên giá trị của tứ phân vị thứ hai là ${Q_2=x_{18}=3}$.\\
		Tứ phân vị thứ nhất là trung vị của mẫu: ${x_1 ; x_2 ; \ldots ; x_{17}}$. Do đó: ${Q_1=x_9=2}$.\\
		Tứ phân vị thứ ba là trung vị của mẫu: ${x_{19} ; x_{20} ; \ldots ; x_{35}}$. Do đó: ${Q_3=x_{27}=4}$.\\
		Vậy $T=Q_{1}+Q_{2}+Q_{3}=9$.
	}
\end{ex}
\Closesolutionfile{ans}
\indapan{6}{ans/ans-KQ-C5-B3-D2}